% !TEX root = ../tecnologia-en-marcha.tex
% !TEX encoding = UTF-8 Unicode
% =============================================================================
% Marcas de revisión del borrador.
%
%   \pendiente{...}  algo que falta escribir o decidir (rojo)
%   \revisar{...}    una afirmación del texto que hay que verificar (naranja)
%
% Antes del envío, cambiar \notastrue por \notasfalse: las marcas desaparecen
% del PDF sin tocar el texto.
% =============================================================================
\newif\ifnotas
\notastrue

\newcommand{\pendiente}[1]{%
  \ifnotas{\color{red!70!black}\textbf{[Pendiente:} #1\textbf{]}}\fi}
\newcommand{\revisar}[1]{%
  \ifnotas{\color{orange!85!black}\textbf{[Revisar:} #1\textbf{]}}\fi}

% "Fuente:" debajo de cuadros, como pide la plantilla.
\newcommand{\fuente}[1]{\par\smallskip{\centering\small Fuente: #1\par}}
